%%%%%%%%%%%%%%%%%%%%%%%%%%%%%%%%%%%%%%%%%%%%%%%%%%%%%%%%%%%%%%
%%     ABSTRACT IN ITALIAN LANGUAGE AND ACKNOWLEDGMENTS     %%
%%%%%%%%%%%%%%%%%%%%%%%%%%%%%%%%%%%%%%%%%%%%%%%%%%%%%%%%%%%%%%
\cleardoublepage

%-----------------------------------------------------------------------------
% SOMMARIO
%-----------------------------------------------------------------------------
\section*{Abstract in lingua italiana}
I Large Language Model (LLM) alimentano sempre più spesso interfacce in linguaggio naturale
verso dati strutturati, ma il loro impiego in produzione solleva questioni che vanno oltre
l'accuratezza, ovvero la latenza di inferenza e il costo energetico. Per studiare questi
aspetti, la tesi realizza un \textit{Sales Agent}: un sistema basato su LLM, orchestrato con
LangGraph, che traduce una domanda di business in una query SQL, ne interpreta i risultati in
forma testuale e, opzionalmente, genera un grafico. Una
configurazione gerarchica in YAML espone i parametri di inferenza studiati
(parametri di campionamento, best-of-$N$, budget di token e profondità del raffinamento
Chain-of-Thought). Il contributo
principale è un profiling sistematico che misura congiuntamente la qualità dell'output, il tempo
di esecuzione e il consumo energetico per singolo step (tramite CodeCarbon). La qualità è
valutata su 20 domande analitiche, stratificate su quattro livelli di difficoltà SQL, lungo tre
dimensioni: un punteggio graduato di similarità tabellare rispetto alla ground truth per il recupero dei dati e due
punteggi LLM-as-judge per l'analisi testuale e la visualizzazione. Cinque esperimenti a step
isolato definiscono una baseline per tre modelli open-weight (Gemma~4~E4B, Gemma~4~26B,
Mistral~Small~3.2~24B) e poi variano uno step dell'agente alla volta, concludendo con una
conferma di Pareto. Sono inoltre progettati per verificare un'assunzione comune nel deployment
locale di LLM, ovvero che i modelli thinking più grandi siano necessariamente la soluzione
migliore. I risultati mostrano invece che la sola scelta del modello non spiega le prestazioni: gli
effetti dominanti emergono dall'interazione tra famiglia del modello, step dell'agente,
difficoltà del prompt e impostazioni di inferenza. Mistral è il modello non ottimizzato più
efficiente ma crolla sui prompt di aggregazione più difficili; Gemma~26B è altamente
ottimizzabile ma energeticamente costoso; Gemma~E4B, corretto con raffinamento Chain-of-Thought
sullo step di lookup a temperatura più bassa, raggiunge la qualità strict più alta e affianca
Mistral, a basso consumo, sulla frontiera di Pareto finale. L'analisi di sensitività mostra
inoltre che gli effetti degli iperparametri sono fortemente specifici per step, che budget di
token più ampi aiutano solo le coppie modello-step realmente troncate, e che il raffinamento
Chain-of-Thought conviene solo quando il suo feedback corrisponde al tipo di errore: è
l'intervento più efficace per il lookup SQL verificabile, ma dannoso per la visualizzazione, dove
gli errori sono semantici e non eseguibili. La tesi distilla tre raccomandazioni per agenti dati
LLM attenti all'energia: ottimizzare gli iperparametri per singolo step dell'agente anziché
globalmente, allocare il compute a tempo di inferenza in base al tipo di errore, e riportare
accuratezza ed energia per livello di difficoltà del prompt, poiché il modello migliore per i
compiti semplici non è il migliore per quelli difficili.
\vspace{15pt}
\begin{tcolorbox}[arc=0pt, boxrule=0pt, colback=bluePoli!60, width=\textwidth, colupper=white]
    \textbf{Parole chiave:} LLM data agents, Energy-aware profiling, Quality-energy trade-off, Local open-weight LLMs, LLM-as-judge evaluation
\end{tcolorbox}
